\section{Finite Cohen Forcing}

The continuum reading at the end of the previous section cannot be accepted
as a gift. A Cauchy process needs finite conditions that make its approach
honest. Finite Cohen forcing is the chapter's name for that discipline. It
does not import an unbounded set-theoretic machine. It names a finite,
\(\epsilon\)-bounded grammar of spline atoms, supports, extensions,
compatibility tests, and dense decisions. The continuum reading is forced
because every admissible finite condition can be extended until the relevant
question is decided within the bound.

A spline atom is the smallest local completion unit the chapter needs. It
has an anchor, a bounded region of support, and a value or activity reading
that can be compared with neighboring atoms. An atom by itself is not a
continuum. It is a finite obligation: here is a local piece of the ledger,
here is what it carries, here is the bound under which it may be extended.
A spline condition is a finite collection of such atoms arranged so that
their overlaps are compatible.

Finite support is essential. If a condition were allowed to carry infinitely
many active obligations at once, it would no longer be a measurement record.
It would pretend to have spent all future refinements in advance. Finite
support says that only finitely many atoms are currently active. Other atoms
may be added later by extension. The grammar therefore preserves the
halting horizon even while it points toward completion.

Extension is the order of forcing. A stronger condition does not abandon the
weaker one. It preserves the commitments already made and adds refinements,
smaller bounds, or more local decisions. If a proposed extension contradicts
an earlier atom on their shared support, it is not an extension at all. It is
a different ledger. Compatibility is the grammar's refusal to let future
precision erase past record without paying for the revision.

Dense decision is the forcing heart. For every question the grammar is
authorized to ask, and for every current condition that has not answered it,
there must be an admissible extension that decides it up to the current
\(\epsilon\). The word dense means that undecided conditions are not allowed
to sit in isolated pockets forever. Wherever the finite ledger stands, the
question can be pressed further by an extension. The decision may be
positive, negative, or bounded as unreadable, but it must be made in the
language of the finite condition.

This is how the continuum is forced rather than assumed. The ledger does not
survey all possible points. It builds a chain of finite conditions such that
every required local question is eventually decided within a shrinking
bound. Two completed readings are identified when no finite dense decision
can separate them. A proposed continuum distinction that cannot be recovered
by any such finite extension is forbidden as unearned structure.

The Cantor--G\"odel--Cohen example serves as a warning. Different
completions may agree on the finite records a grammar has chosen to carry.
Consistency alone is therefore too weak. A completion is admissible here
only when its distinctions are recoverable from finite refinement. Finite
Cohen forcing gives the grammar a way to say which completions are forced by
the ledger and which are optional stories placed above it.

The compact-disc threshold example gives the same warning in concrete form.
A coarser reader may record fewer events than a finer reader. When the
records are merged, the finer distinctions force refinements on the coarser
history where compatibility requires them. But distinctions below every
available threshold do not enter the record merely because a continuum could
name them. They must be forced by an admissible extension of the ledger.

For Chapter 04, the forced continuum reading is the stable home of the
pair-production grammar. The electron sign, the tilted positron sign, the
boundary-generated gauge, the split, and the zero-mesh relation all require
finite decisions to survive refinement. A condition may decide a local sign.
Another may decide a boundary compatibility. Another may shrink a geometric
residual. The forcing process keeps extending until every required decision
is dense in the refinement order.

This does not erase finitude. It honors it. The continuum reading is the
quotient of refinement histories that finite dense decisions can no longer
distinguish. The grammar carries atoms, bounds, supports, and extensions. It
forbids unrecorded intermediate structure. It identifies compatible
late-stage readings. It forces the continuum as the only reading left after
all finite questions have been pressed under shrinking bounds.

The chapter has therefore reached a full circle. It began with the concrete
second variation and read the electron. It tilted the baseline and read the
positron. It isolated the two grants that make those readings lawful. It let
boundary geometry generate the gauge. It split geometry and gauge to forbid
interior mixing. It read the electromagnetic face as the Cauchy limit of the
geometric face. Finally, it forced the continuum reading by finite
conditions. The next chapter inherits the consequence: when the interior
mixed term is forbidden and the continuum reading is forced, radiation must
be read at the boundary.

